\begin{figure}[ht]
\centering
\begin{tikzpicture}[x=1cm,y=1cm,>=Latex,font=\small]
\draw[->,gray] (-.4,0)--(9,0) node[right] {$x$};
\draw[->,gray] (-.4,0)--(-.4,2.9) node[above] {$z$};
\draw[very thick,blue!65!black] (.7,0)--(2.3,0);
\draw[very thick,green!45!black] (5.7,0)--(7.3,0);
\node[below] at (1.5,0) {$C_i\times\{0\}$};
\node[below] at (6.5,0) {$\Theta_i(C_i)\times\{0\}$};
\draw[->,thick,blue!65!black] (1.5,.15)--(1.5,2.2)
 node[midway,left] {lift};
\draw[->,thick] (1.7,2.35)--(6.3,2.35)
 node[midway,above] {reciprocal affine interpolation};
\draw[->,thick,green!45!black] (6.5,2.2)--(6.5,.15)
 node[midway,right] {lower};
\draw[dashed,gray] (-.25,2.35)--(1.25,2.35);
\node[left] at (-.4,2.35) {$h_i$};
\end{tikzpicture}
\caption{The three scalar-potential phases for one instruction. This is a schematic projection onto the $x,z$ coordinates; the second planar coordinate is suppressed. Distinct branches use separate heights. Source rectangles are disjoint among themselves, as are image rectangles, while a source may overlap an image. During the middle phase the first coordinate stays between its endpoints.}
\label{fig:scalar-stages}
\end{figure}
